\section{The Emperor}

\begin{quotex}
The ones who truly love their traditions don't take them too seriously. They march to get their heads shot off with a joke on their lips. And the reason is that they know they're going to die for something intangible, something sprung from their fancy, half humor, half humbug. Or perhaps it's a little more subtle. Perhaps hidden away in their fancy is that pride of the blueblood, who refuses to look foolish by fighting for an idea, and so he cloaks it with bugle calls that tug at the heart, with empty mottoes and useless gold trim, and allows himself the supreme delight of giving his life for an utter masquerade. That's something the Left has never understood, and that's why its contempt is so heavy with hate. When it spits on the flag, or tries to piss out the eternal flame, when it hoots at the old farts loping by in their berets, or yells ``Women's Lib!'' outside the church, at an old-fashioned wedding (to cite just some basic examples), it does so in such a grim, serious manner --- like such ``pompous assholes,'' as the Left would put it, if only it could judge. The true Right is never so grim. That's why the Left hates its guts, the way a hangman must hate the victim who laughs and jokes on his way to the gallows. The Left is a conflagration. It devours and consumes in deadly dull earnest. (Even its revels, appearances notwithstanding, are as grisly an affair as one of those puppet parades out of Peking or Nuremberg.)The Right is different. It's a flickering flame, a will-o'-the-wisp in the petrified forest, flitting through the darkness...

\flright{\textsc{Jean Raspail}, \emph{The Camp of the Saints}\footnote{\url{https://en.wikipedia.org/wiki/Jean_Raspail} (1925--) a French author, and explorer.}}
\end{quotex}


The Number of the Emperor is Four; this is the sign of Earth's four elements, and also the procession of the movement of Being out of the perfect triangle-circle of One-Two-Three, into the world of material creation, the union of spirit and matter. The Emperor has authority because he has achieved existence, he knows something, \& he is capable of acting (1-2-3): he can therefore assume the post of emperor because he is ``earth of earth'', ``the most human of humans'', King David -- a shepherd, a hero, a champion, a warrior, an outlaw, a king, a sinner, a saint. He is the manifestation of the Divine in the form of he who bears the sceptre, and is (therefore) the wielder of the sword not in vain.

The rise of the emperor to the post appointed for him by God occurs in the topsy-turvy world: for this reason, the emperor is ``out on the green grass'', under the skies. His authority originated in the fact that the power of his being is such that he can hold court in the wilderness, and indeed (most likely) this may be where he plants the seeds that lead to his rule. As Rosenstock-Huessy argues in \emph{Out of Revolution}, real political authority is worth nothing unless it is so powerful that it first begins ``under the stars'' or the clear blue skies. We reject the ``whiff of Hegelianism'' that wafts from Huessy's\footnote{\url{http://plato.stanford.edu/entries/rosenstock-huessy/}} work; however, he is correct to say that political authority changes with the revolution of the stars, and that even the secular revolutions (which aped legitimate procession of government) adhere to this change that which falls from the skies first (particularly when government is corrupt and asleep at the wheel). \emph{Auctoritas} comes from the one who is self-restrained and crucified: that is, the emperor does not need the outer symbols to rule -- they come to him like gravity draws smaller objects: he is comfortable at first, not with worldy power, but with naked Being. Like the sage on the mountaintop, the emperor sits by the river banks, or in the woody dells, or out in a cow pasture. He dwells outside the ``city'', when the city is falling apart and degenerate. He is at home within his own skin, and is free from worldly A influences\footnote{\url{http://unurthed.com/2009/10/21/mouravieffs-correction/}} which dominate the ``kingdom of this world''. He is King Alfred living in the swamps, awaiting his chance to reconquer England from the Vikings.

This is the opposite of the career politician, the demagogue, or the opportunist. His power comes from power over self, not over other people. He holds sway over others, and carries or embodies real authority over them, and as a result, he exercises legitimate power. Tomberg makes the strong case that compulsion is superfluous for those with real authority: ``where there is authority, there is present the breath of sacred magic filled by the rays of light of gnosis emanated from the profound fire of mysticism, there compulsion is superfluous''. Viewed this way, Leftism or revolutionary thought is basically like a mental disease, which is healed by true action on the body politic through the esoteric organs. This is consistent with the doctrine of the Gnostic Church as the ``heart'' of the organized Petrine and Pauline Churches.

``The Emperor has renounced movement by means of his legs (they are crossed) and action by means of his arms (one hitches the belt, the other holds the sceptre).'' The ``instinctive and impulsive'' nature of the Emperor is restrained, by self-restraint acting under Divine influence. It is this, and this alone, which allows him to hold the post of guardian of the office of Empire. When King Alfred made peace with the Danes, but had Guthrum baptized, he secured the Anglo-Saxon kingdom of England and the Danelaw as a united realm, bringing peace to the land. He could have hid vengeance in his heart, but instead, he converted his enemy, quite literally, to a friend. It is this type of action which the Holy Roman Emperor aimed at, being king of kings and lord of lords, the means by which the petty rivalries of monarchs (which would later destroy Europe in repeated conflagrations) were brought to heel before the sacred body of the King\footnote{\url{http://www.amazon.com/The-Kings-Bodies-Ernst-Kantorowicz/product-reviews/0691017042}}. Tomberg remarks that the Middle Ages contain a wealth of political theory (which is unique, even from classical Greece and Rome) on the post of the Emperor: under this list, we include Dante (who sided with the White Guelphs\footnote{\url{http://www.youtube.com/watch?v=UcGDRYQd-uQ&feature=related}} in the Ghibelline controversy) and John of Paris\footnote{\url{http://en.wikipedia.org/wiki/John_of_Paris}}.

``The emperor has renounced personal opinion to receive revelation of Truth, personal action in order to become a vehicle of sacred magic, the way of personal development in order to be guided by the Master of the Way, and his own personal mission in order to be charged with a mission from above.'' This makes him the ``source of Law and Order''.

I think I have quoted Bonaventura before\footnote{\url{http://noxrpm.com/post/2609941502/non-coerceri-maximo-contineri-minimo-divinum}}: That which is not confined by the Great, but is contained by even the smallest, it is that which is Divine. God (the Self beyond Self) becomes King in respect of those who freely worship Him, and He is crucified in relation to those who choose to fight or struggle with that freedom against the Font of Freedom. In both cases, a withdrawal is accomplished in order to establish a realm of freedom. Hitler and Napoleon, both, attempted to occupy the post of the shadow of the Emperor, but they took up the sword to do it. That is, they placed their trust in worldly machinations: power and the grapeshot or Panzers that were seen to produce such.

There is an Emperor of Europe, of Christendom. The post is occupied. It in our day is occult, in shadow, but God does not allow a vacuum in Nature: the man who has a vacuum of Ego and Pride will be courted by the Spirit for the post of Emperor, and there is undoubtedly an occulted Emperor today, in 2014. The complete synthesis of mysticism-gnosis-magic which is effected in the Emperor is also, Tomberg states, the definition of \textbf{\emph{initiation}}, where eternity and the present moment meet together.

Mysticism becomes true without falsehood (Gnosis), and most certain (magic), and then absolutely certain in the light of pure thought as both subjective and objective experience (Hermetic philosophy): the Emperor is the hermetic synthesis of the emanations of God's light upon the world, though he may be unappreciated, unknown, hidden, or occulted.

``Presumption? It would be a monstrous presumption if it were a matter of human invention instead of revelation from above.''

``Ask'' (mysticism or touch), ``seek'' (gnosis or hearing), ``knock'' (magic or sight/vision) -- Luke 11:9. These become \emph{comprehension}. This is \emph{initiation}. The goal of initiation is \emph{depth}, or ``niveau''. For the post of emperor, the aim is the realization of patriarchy, the ``father-man'', the most human of all men, the heir of David. It is the four wounds spoken of which mark the Emperor, and provide him with the concrete experience to fulfill his destiny, to be anonymous, or (rather) synonymous with the post of Emperor itself : hence it is said, King Arthur only sleeps, as does the ``King under the Mountain'' (the Holy Roman Emperor). If the king sleeps, he shall return.

Indeed he is already here.

This bears reflection upon in the present day, for those who want to rebuild the monarchies. It was the Empire which gave birth to kings, the Holy Empire, and not the kings who built up the Empire. Rather than seeking to restore a particular pedigree or bloodline, thus exposing these men and the movement to an undoubted reaction which will be most violent, more thought should be put into finding the occulted Emperor, and also, collating and honing the methods and practices which will produce divine candidates of this quality. There is a lot of synthesis to be done in esoteric studies, and it can only be done rightly by those who have experienced it. The Emperor will require men of depth around him to respond to his authority, \& also to embody emperorship in their little field of influence. Has anyone given thought to joining the Gnosis seminar? Mouravieff gives a practical method for recovering the initiation of Christendom.

One ought to try to recognize Emperor within one's self, by seeking that which is truly Just. As Cologero has outlined, justice is the synthesis of moderation (the man of the senses), courage (the man of the emotions) and wisdom (the man of the intellect). You can see the same symbolism of numeric progression: what does Courage have to do with Gnosis? This is surprising, as one would expect Intellect to correlate with Gnosis, and Courage with Magic. When such a deliberate contradiction to our ``everyday'' wisdom or literalness is generated, there is an opportunity to look more deeply. The man of Gnosis needs courage as the primary virtue because it goes so strongly against the ``naive realism'' which dominates worldly thinking. The man of Magic requires intellect, because he has already mastered courage (or courage has mastered him) and he requires mastery of the science and art, since Magic is potentially quite dangerous. Fast (senses), Watch (master passions), and Pray (effect the use of Magic under the name of God alone).


\flrightit{Posted on 2014-12-26 by Logres}

\begin{center}* * *\end{center}

\begin{footnotesize}
\begin{sffamily}
\texttt{Olavsson on 2014-12-27 at 14:28 said:}

Most beautifully written and absolutely on target. The coming of the True Emperor is certain, and the realization of that truth is enough to keep the fire burning no matter how long we have to wait. Always remain loyal to the end.

\hfill

\texttt{TeenyTinyTarot on 2014-12-28 at 18:18 said:}

\begin{quotex}\footnotesize
What I assert is this,---that a man ought to be in serious earnest about serious things, and not about trifles; and that the object really worthy of all serious and blessed effort is God, while man is contrived, as we said above, to be a plaything of God, and the best part of him is really just that; and thus I say that every man and woman ought to pass through life in accordance with this character, playing at the noblest of pastimes, being otherwise minded than they now are.
\flright{\textsc{Plato}, \textit{Laws}, 803c.\footnote{Just read another translation of this as quoted by Peter Kingsley in \textit{Ancient Philosophy, Mystery, and Magic}, page 170.}}
\end{quotex}

\hfill

\texttt{kenthamaker on 2014-12-29 at 10:35 said:}

This is a beautiful post and a fecund inspiration for a new year. Thank you.

\hfill

\texttt{logres on 2014-12-29 at 13:21 said:}

Thank you for reading. That is a good quote, TTT: as ``Samwise'' on the old Spengler Forum at Asia Times Online used to argue, religion is or ought to be a ``task'' -- that is, God, immortality, and the soul have to be ``achieved'' by somebody. So we can assign ourselves the yoke of virtue, if we are listening to the highest parts of our soul. This was the entire yearning of the old classical tradition.

\hfill

\texttt{Space Lizard on 2014-12-29 at 15:49 said:}

From the ashes a fire shall be woken,

A light from the shadows shall spring;

Renewed shall be blade that was broken,

The crownless again shall be king.

\hfill

\texttt{Logres on 2014-12-31 at 22:33 said:}

But the Consul's brow was sad,

And the Consul's speech was low, 210

And darkly looked he at the wall,

And darkly at the foe;

``Their van will be upon us

Before the bridge goes down;

And if they once may win the bridge, 215

What hope to save the town?''

Then out spake brave Horatius,

The Captain of the gate:

``To every man upon this earth

Death cometh soon or late. 220

And how can man die better

Than facing fearful odds

For the ashes of his fathers

And the temples of his gods,

``And for the tender mother 225

Who dandled him to rest,

And for the wife who nurses

His baby at her breast,

And for the holy maidens

Who feed the eternal flame,--- 230

To save them from false Sextus

That wrought the deed of shame?

\hfill

\texttt{Aquila on 2015-01-15 at 12:30 said:}

According to my information, Dante sided with the Ghibellines. Now, what is White Guelphs? Never heard about that.

\hfill

\texttt{John Fitzgerald on 2021-02-07 at 15:12 said:}

I have been aware of the existence of this site for a number years, but for some reason it's passed me by until this very weekend. It's been the long-overdue eruption of Dante into my life (my fault for the delay, not his) that has brought me here.

I have only just read this powerful reflection -- or evocation, rather -- of the Emperor this evening. Yet this piece I wrote for the Albion Awakening blog in May 2017 seems uncannily similar, certainly with regards to the Emperor appearing (or reappearing) on the periphery as opposed to the centre.

\url{http://albionawakening.blogspot.com/2017/05/the-return-of-constantine.html}


I had read and absorbed Tomberg though, and continue to do so, as it's a life-time's endeavour, of course. So maybe that accounts for the similarities.

Thank you and God bless,

John Fitzgerald

Cymru / Manchester

\end{sffamily}
\end{footnotesize}

